\documentclass[10pt,a4paper]{amsart}
\usepackage[margin=19mm]{geometry}
\usepackage{amsmath,amssymb,mathtools}
\usepackage[scr=rsfs,cal=euler]{mathalfa}
\usepackage{xfrac}
\usepackage[unicode,hidelinks,bookmarks=false]{hyperref}
\newcommand{\ie}{i.e.\ }
\newcommand{\cf}{cf.\ }
\newcommand{\resp}{respectively\ }
\newcommand{\Subsection}{\subsection}
\providecommand{\nobreakdash}{\nobreak\hbox{-}}
\setlength{\emergencystretch}{2em}
\multlinegap0pt
\usepackage{tikz}
\usepackage{amssymb}
\usepackage{float}
\usepackage{bm}
\usepackage{xcolor}
\usepackage{enumerate}
\usepackage[shortlabels]{enumitem}
\usepackage[normalem]{ulem}
\newtheorem{lemma}{Lemma}
\title{The Multinomial Allocation Model and the Random Box Load}
\author{Serik Sagitov}
\date{Source: arXiv:2604.15152v2 (CC BY 4.0)}
\begin{document}
\maketitle

\noindent\emph{Source and licence.} The Multinomial Allocation Model and the Random Box Load. Version arXiv:2604.15152v2 (CC BY 4.0), distributed by arXiv under the Creative Commons Attribution 4.0 International licence, http://creativecommons.org/licenses/by/4.0/, as recorded in the arXiv licence metadata for this exact version and shown on the abstract page https://arxiv.org/abs/2604.15152v2. Full licence notice: \texttt{licenses/arxiv-2604.15152-CC-BY-4.0.txt}.\par\smallskip

\noindent\textit{Editorial scope.} Counted unit: the article's lemma L5 (source0.tex lines 271--287). The article's complete original proof of it is delivered with it (source0.tex lines 289--358, one \texttt{proof} environment of 70 lines). No mathematics is written, repaired or re-derived: the statement and the proof are the article's own lines, in the article's own notation, and no retained line is substituted or re-flowed.  No further source-local context is needed: every cross-reference of every retained line resolves inside the two delivered spans.  Nothing was omitted from any retained span, so the package carries no omission record.  No retained line contains a citation, so the wrapper carries no bibliography and no bibliography entry.  No retained line contains a figure environment, an image-inclusion command or drawing code, so no image is delivered and the package declares no asset dependency.  Editorial formatting only, in addition: an AMS \texttt{amsart} preamble that restates the article's own declarations and macros used by the retained lines, namely the environments \texttt{lemma} on separate counters, together with the standard packages those lines need.  The retained environments print in this document as Lemma 1 in order of appearance, whereas the article prints the same items with the numbering of its own full document.  The complete native source of the article as distributed in the arXiv e-print for this exact version is bundled unchanged as \texttt{sources/198583/source0.tex}.
\begin{lemma}\label{L5}
For $n \ge 1$ and $0\le r\le n$, define $\delta_n(r,x)$ by
\[
\left(1-\frac{x}{n}\right)^{n-r} = e^{-x} + n^{-1}\delta_n(r,x).
\]
If $0 \le x \le \frac{n}{2}$, then
\begin{equation}\label{del}
 -2 \le \delta_n(r,x) \le  r2^{r-1},
\end{equation}
and moreover, $\delta^*_n(r,x)$ defined by
\[\delta_n(r,x)=e^{-x}x(r-x/2)+n^{-1}\delta^*_n(r,x),
\]
satisfies 
\begin{equation}\label{del*}
-4-r\le \delta^*_n(r,x)\le x^2r(r+1) 2^{\,r+1}.
\end{equation}
\end{lemma}

\begin{proof}
We use the representation
\[
\left(1-\frac{x}{n}\right)^{n-r}
= \left(1-\frac{x}{n}\right)^{-r} e^{\,n\ln(1-x/n)}
\]
and repeatedly apply Taylor's theorem with Lagrange remainder. By
\[
\ln(1-y) = -y - \frac{y^2}{2(1-\theta y)^2}, \qquad \theta \in [0,1],
\]
we have
\[
-y - 2y^2\le \ln(1-y)\le -y,\quad 0\le y\le 1/2.
\]
With $y = \frac{x}{n}$, it follows that
\[
-x - \frac{2x^2}{n}
\le n\ln\left(1-\frac{x}{n}\right)
\le -x,\quad 0\le x\le n/2.
\]
Exponentiating and using the inequality $e^{-t} \ge 1 - t$ for $t \ge 0$, we obtain
\[
e^{-x}\left(1 - \frac{2x^2}{n}\right)
\le e^{\,n\ln(1-x/n)}
\le e^{-x}.
\]

On the other hand,\[
1 \le \left(1-\frac{x}{n}\right)^{-r}
\le 1 + \frac{xr}{n}2^{\,r}.
\]
Combining the bounds, we find
\[
-2x^2 e^{-x}
\le \delta_n(r,x)
\le e^{-x} x r 2^{\,r},
\]
which  yields \eqref{del} in view of $e^{-x}x<1/2$ and $e^{-x}x^2<1$.

The higher order Taylor's expansion,
\[
\ln(1-y) = -y - \frac{y^2}{2}- \frac{y^3}{3(1-\theta y)^3}, \qquad \theta \in [0,1],
\]
yields
\[
-x - \frac{x^2}{2n}- \frac{8x^3}{3n^2}
\le n\ln\left(1-\frac{x}{n}\right)
\le -x- \frac{x^2}{2n},\quad 0\le x\le n/2.
\]
Exponentiating and using 
\[e^{-t} = 1 - t+\frac{t^2}{2}e^{-\theta't},\quad t \ge 0, \qquad \theta' \in [0,1],\]
we obtain
\[
e^{-x}\left(1 - \frac{x^2}{2n}- \frac{8x^3}{3n^2}\right)
\le e^{\,n\ln(1-x/n)}
\le e^{-x}\left(1 - \frac{x^2}{2n}\right)+\frac{x^4}{4n^2}.
\]

On the other hand,
\[
 \left(1-\frac{x}{n}\right)^{-r}
= 1 + \frac{xr}{n}+\frac{r(1+r)x^2}{2n^2}\left(1-\frac{\theta''x}{n}\right)^{-r-2}, \qquad \theta'' \in [0,1],
\]
yields
\[
1 + \frac{rx}{n}\le \left(1-\frac{x}{n}\right)^{-r}
\le 1 + \frac{rx}{n}+\frac{r(1+r)x^2}{n^2}2^{r+1},
\]
Combining these bounds, and using $e^{-x}x^3<3/2$, we establish \eqref{del*}.
\end{proof}

\end{document}
